\documentclass[11pt]{article}
\usepackage[margin=1in]{geometry}
\usepackage{amsmath,amssymb,amsthm,mathtools}
\usepackage[T1]{fontenc}
\usepackage{lmodern,microtype}
\setlength{\emergencystretch}{2em}
\usepackage[hidelinks]{hyperref}
\newtheorem{theorem}{Theorem}[section]
\newtheorem{lemma}[theorem]{Lemma}
\newtheorem{proposition}[theorem]{Proposition}
\newcommand{\HH}{\mathsf H}
\newcommand{\Int}{\operatorname{Int}}
\newcommand{\Asc}{\operatorname{Asc}}
\title{A fully supported local intricacy maximum\newline among stationary binary one-step Markov laws}
\author{Research manuscript for MathDB333569\\\small Draft for independent paper review; no priority assertion}
\date{1 November 2026}
\begin{document}
\sloppy
\maketitle
\begin{abstract}
For the coordinate partition of the two-sided full binary shift and uniform weights on finite subsets, we prove existence of a fully supported stationary one-step Markov law that locally maximizes intricacy relative to all stationary one-step Markov laws, with their relative weak topology. We derive the limiting entropy series directly from finite selected-coordinate entropies. A compactness argument gives a maximum on the symmetric line. The decisive calculation proves strictly negative transverse second derivative at every positively correlated symmetric law; uniform geometric derivative bounds then promote the diagonal maximum to a two-dimensional local maximum. Continuous recovery of transition probabilities from clopen cylinder masses covers every nearby comparison law, including the full class's potentially degenerate members. An accompanying unchanged Lean certificate proves the actual-law endpoint and convergence of the finite expressions at the witness and nearby laws. Formal correctness and the bounded literature comparison are distinct claims.
\end{abstract}

\section{The question and the exact assertion}
Let $X=\{0,1\}^{\mathbb Z}$ with its product topology and Borel sigma algebra, and let $T x$ have coordinates $(Tx)_i=x_{i+1}$. Let $\alpha$ be the coordinate-zero partition. For a probability law $\mu$ and finite $S\subset\mathbb Z$, write
\[
 \HH_\mu(S)=-\sum_{a\in\{0,1\}^{S}}\mu\{x:x|_S=a\}\log\mu\{x:x|_S=a\}.
\]
All logarithms are natural, $0\log0=0$, and $\HH_\mu(\varnothing)=0$. Put $B_N=\{0,\ldots,N-1\}$ for $N\geq1$. The finite uniform-coefficient quantities are
\begin{align}
 \Asc_N(\mu)&=\frac1{N2^N}\sum_{S\subset B_N}\HH_\mu(S),\label{asc}\\
 \Int_N(\mu)&=\frac1{N2^N}\sum_{S\subset B_N}
 [\HH_\mu(S)+\HH_\mu(B_N\setminus S)-\HH_\mu(B_N)]\notag\\
 &=2\Asc_N(\mu)-\frac1N\HH_\mu(B_N).\label{finite}
\end{align}
The last identity follows because complementation permutes the subsets. Whenever this sequence converges, denote its limit by $\Int(\mu)$. This is intricacy for the rank-zero partition and uniform subset weights, not a supremum over partitions.

A stationary one-step Markov law means a probability law with some stochastic matrix $P$ on $\{0,1\}$ and stationary probability vector $p$, whose every translated block satisfies, for $k\in\mathbb Z$ and $n\geq0$,
\begin{equation}
 \mu[x_k=a_0,\ldots,x_{k+n}=a_n]
 =p_{a_0}\prod_{j=0}^{n-1}P_{a_j a_{j+1}}.\label{block}
\end{equation}
Let $\mathcal M_1$ be the class of all such laws. In this definition $P$ may have zero entries, $p$ may have zero entries, and the law need not be ergodic. Give $\mathcal M_1$ the relative topology inherited from weak convergence of probability measures on $X$.

\begin{theorem}\label{main}
There exist $\mu_*\in\mathcal M_1$ and a weak neighborhood $U$ of $\mu_*$ such that:
\begin{enumerate}
\item $\mu_*$ is invariant under every integer shift and gives positive mass to every nonempty open subset of $X$;
\item $\Int_N(\mu_*)$ converges, and $\Int_N(\nu)$ converges for every $\nu\in U\cap\mathcal M_1$;
\item $\Int(\nu)\leq\Int(\mu_*)$ for every $\nu\in U\cap\mathcal M_1$.
\end{enumerate}
\end{theorem}
This resolves the existence assertion in Petersen's Conjecture 4.6.8 in arXiv:1607.02425v1 \cite{P}, and Petersen--Wilson's Conjecture 6.3 in arXiv:1512.01143v6 \cite{PW}, in the cached versions specified. It proves neither global optimality nor uniqueness, no numerical value of a maximizing parameter, and no assertion about higher-order Markov laws or all invariant laws.

\section{A stationary chart and its two-sided laws}
For $0<l<1$, $|t|<1$, set
\begin{equation}
 \pi_t=\left(\frac{1+t}{2},\frac{1-t}{2}\right),\qquad
 P(l,t)=\begin{pmatrix}
 1-(1-l)(1-t)/2 &(1-l)(1-t)/2\\
 (1-l)(1+t)/2&1-(1-l)(1+t)/2
 \end{pmatrix}.\label{chart}
\end{equation}
Both coordinates of $\pi_t$ and all four entries of $P(l,t)$ are positive. Indeed the off-diagonal entries lie strictly between zero and one, since $0<1-l<1$ and $0<1\pm t<2$. The rows sum to one. If $J_t$ has both rows equal to $\pi_t$, then
\[
 P(l,t)=lI+(1-l)J_t,\quad J_t^2=J_t,\quad \pi_tJ_t=\pi_t.
\]
It follows that $\pi_tP(l,t)=\pi_t$ and, for integers $i\geq1$,
\begin{equation}
 P(l,t)^i=J_t+l^i(I-J_t)=P(l^i,t).\label{powers}
\end{equation}

For completeness these data specify an actual two-sided probability law. Assign every interval $[k,k+n]$ the mass in \eqref{block}, with $p=\pi_t$ and $P=P(l,t)$. Summing out a right endpoint uses row sums one; summing out a left endpoint uses stationarity. Iterating gives compatible marginals on arbitrary finite coordinate sets (embed each in an interval and sum out the complement). They are nonnegative normalized probabilities. The countable product extension theorem gives a Borel probability measure $\mu_{l,t}$ on $X$, unique because finite-coordinate cylinders generate its Borel sigma algebra and form a pi-system. The assignment is independent of $k$, so uniqueness proves invariance under all integer shifts. Every full interval word has strictly positive mass; any nonempty finite-coordinate cylinder contains an interval word, and every nonempty open set contains a finite-coordinate cylinder. Thus $\mu_{l,t}$ has full open support.

The unchanged Lean proof constructs this measure through two one-sided trajectory laws conditioned at zero and glues them. This equivalent construction uses detailed balance, valid here since
\[
 (\pi_t)_0P_{01}=(\pi_t)_1P_{10}=(1-l)(1-t^2)/4.
\]
Reversing paths then gives the chronological block probabilities \eqref{block}, including intervals crossing zero; uniqueness identifies the resulting law with the extension just described.

Writing $u=P_{00}$, $v=P_{11}$, the inverse chart is
\begin{equation}
 l=u+v-1,\qquad t=\frac{u-v}{2-u-v}.\label{inverse}
\end{equation}
Indeed $u+v=1+l$, $u-v=(1-l)t$. The denominator equals $1-l>0$. Conversely substitution of \eqref{inverse} into \eqref{chart} returns $(u,v)$ whenever $2-u-v\ne0$. Around each symmetric point $u=v=(1+l)/2$ with $0<l<1$, these maps are continuous inverses between open neighborhoods. In particular the chart contains every matrix in a sufficiently small neighborhood of that point, not just a prescribed symmetric slice. For an interior matrix the unique stationary probability vector is
\begin{equation}
 p_0=\frac{1-v}{2-u-v},\qquad p_1=\frac{1-u}{2-u-v}.\label{stationary}
\end{equation}
To verify uniqueness, the stationarity equation gives $p_0(1-u)=p_1(1-v)$; substituting $p_1=1-p_0$ yields \eqref{stationary}. The positive denominator justifies division.

\section{Finite selected entropies and the source limit}
Write
\[
 h(x)=-x\log x-(1-x)\log(1-x),\qquad 0\leq x\leq1,
\]
with continuous endpoint values. This function lies in $[0,\log2]$. Define
\begin{equation}
 H(r,t)=\frac{1+t}{2}h\left(\frac{(1-r)(1-t)}2\right)
       +\frac{1-t}{2}h\left(\frac{(1-r)(1+t)}2\right),\label{H}
\end{equation}
for $0\leq r\leq1$, $|t|\leq1$. The row entropies in $P(r,t)$ have these arguments, using $h(x)=h(1-x)$. Thus $H(l,t)$ is the stationary one-step conditional entropy, and $H(l^i,t)$ is the conditional entropy across gap $i$ by \eqref{powers}. Put $H_0=h((1+t)/2)$.

\begin{lemma}\label{selected}
For $S=\{s_0<\cdots<s_d\}$ nonempty,
\[
 \HH_{\mu_{l,t}}(S)=H_0+\sum_{j=0}^{d-1}H(l^{s_{j+1}-s_j},t).
\]
\end{lemma}
\begin{proof}
Embed $S$ in $[s_0,s_d]$ and sum \eqref{block} over every unselected symbol. For a gap $g$, the intervening sum of products of $g$ transition entries is $(P^g)_{ab}$, by matrix multiplication (induction on $g$). Summing each disjoint gap therefore yields the selected atom mass
\[
 q(a_0,\ldots,a_d)=(\pi_t)_{a_0}\prod_{j=0}^{d-1}
 P(l^{s_{j+1}-s_j},t)_{a_j a_{j+1}}.
\]
All factors are positive. Expand $-q\log q$ using the logarithm of the product and sum. For the initial term, summing all later states gives one by row normalization, so its contribution is $H_0$. In transition term $j$, summing later states gives one; summing earlier states leaves $\pi_t$ at time $s_j$, since $\pi_tP^g=\pi_t$ for every $g$. Its contribution is exactly $H(l^{s_{j+1}-s_j},t)$. This proves the identity, including singleton $S$ with empty transition sum. Empty $S$ separately has a single atom of mass one and entropy zero.
\end{proof}

\begin{lemma}\label{average}
For every $N\geq1$,
\begin{equation}
 \frac{2}{2^N}\sum_{S\subset B_N}\HH_{\mu_{l,t}}(S)
 =2(1-2^{-N})H_0+\sum_{i=1}^{N-1}(N-i)2^{-i}H(l^i,t).\label{count}
\end{equation}
\end{lemma}
\begin{proof}
Each of the $2^N-1$ nonempty subsets contributes one $H_0$. For fixed $a<b$ of gap $i=b-a$, the two times are consecutive in the selected order precisely when $a,b$ are selected and the $i-1$ intervening times are omitted. The other $N-i-1$ coordinates are free; hence exactly $2^{N-i-1}$ subsets contribute this particular adjacent pair. There are $N-i$ possible placements of a pair of gap $i$. Sum the identity of Lemma~\ref{selected} over all subsets and multiply by $2/2^N$. This argument uses linearity of the sum, with no independence assumption on overlapping pairs.
\end{proof}
The block entropy is $H_0+(N-1)H(l,t)$ by Lemma~\ref{selected}. Thus \eqref{finite} gives the exact boundary formula
\begin{align}
 \Int_N(\mu_{l,t})={}&\frac{(1-2^{1-N})H_0+H(l,t)}N-H(l,t)\notag\\
 &+\sum_{i=1}^{N-1}(1-i/N)2^{-i}H(l^i,t).\label{boundary}
\end{align}
For $N=1$ this is zero, which also checks the normalization.

\begin{proposition}\label{series}
For every $0<l<1$, $|t|<1$, the finite intricacy converges and
\begin{equation}
 \Int(\mu_{l,t})=F(l,t):=\sum_{i=1}^\infty2^{-i}H(l^i,t)-H(l,t).\label{F}
\end{equation}
In fact $|\Int_N(\mu_{l,t})-F(l,t)|\leq5\log2/N$.
\end{proposition}
\begin{proof}
The series is absolutely convergent since $0\leq H\leq\log2$. The absolute value of the boundary fraction in \eqref{boundary} is at most $2\log2/N$. The difference between the triangular sum and the infinite sum is at most
\[
 \log2\left(\frac1N\sum_{i=1}^{N-1}i2^{-i}
                  +\sum_{i=N}^{\infty}2^{-i}\right)
 \leq\log2\left(2/N+2^{1-N}\right)\leq3\log2/N.
\]
Here $N\leq2^{N-1}$ for $N\geq1$, proved by induction, gives the last inequality. This proves both convergence and the stated bound.
\end{proof}
The series in \eqref{F} agrees with twice the average sample complexity formula of \cite[Proposition 5.10]{PW} minus the entropy rate. The direct derivation above avoids assuming an unproved identification of an infinite expression with the finite source functional. For stationary endpoints the two conditional directions have equal entropy: $\HH(X_0|X_i)=\HH(X_i|X_0)$ follows by subtracting equal marginal entropies from the joint entropy.

\section{An interior maximum on the symmetric line}
Let $f(l)=F(l,0)$. Formula \eqref{H} gives $H(r,0)=h((1-r)/2)$, and hence
\begin{equation}
 f(l)=\sum_{i=1}^{\infty}2^{-i}
 \left[h((1-l^i)/2)-h((1-l)/2)\right].\label{diagonal}
\end{equation}
This expression extends continuously to $[0,1]$: each term is continuous, bounded in absolute value by $2^{-i}\log2$, and the geometric majorant makes the convergence uniform. At both endpoints all differences vanish, so $f(0)=f(1)=0$. For $0<l<1$, $l^i\leq l$, strictly when $i\geq2$. Since $h'(x)=\log((1-x)/x)>0$ for $0<x<1/2$, each difference in \eqref{diagonal} is nonnegative, and the $i=2$ difference is positive. Therefore $f(l)>0$ throughout $(0,1)$. Compactness yields a maximizer $l_0\in[0,1]$, and positivity and zero endpoint values force
\begin{equation}
 0<l_0<1,\qquad f(l)\leq f(l_0)\quad(0\leq l\leq1).\label{maxdiag}
\end{equation}
No uniqueness or strict longitudinal curvature is required. A diagonal maximum by itself does not prove Theorem~\ref{main}; it remains to control departures from symmetry.

\section{The transverse coefficient and its sign}
Put $x=(1-r)/2$ and $u=x(1-t)$, $v=x(1+t)$. On $0\leq r<1$, $|t|<1$, both $u,v$ lie in $(0,1)$. In the following derivatives $r$ is fixed. Differentiating the weights as well as their arguments in \eqref{H} gives
\begin{align}
 H_t(r,t)&=\tfrac12[h(u)-h(v)]
 +\tfrac{x}{2}[(1-t)h'(v)-(1+t)h'(u)],\label{first}\\
 H_{tt}(r,t)&=-x[h'(u)+h'(v)]
 +\tfrac{x^2}{2}[(1+t)h''(u)+(1-t)h''(v)].\label{second}
\end{align}
Here $h''(x)=-1/x-1/(1-x)$. At $t=0$ these imply
\begin{equation}
 H_{tt}(r,0)=2c(r),\qquad
 c(r)=-\frac{1-r}{2}\log\frac{1+r}{1-r}
       -\frac{1-r}{2(1+r)}.\label{coefficient}
\end{equation}
For example, half of \eqref{second} at zero is
$-xh'(x)+x^2h''(x)/2=-x\log((1-x)/x)-x/(2(1-x))$, giving \eqref{coefficient}. Direct differentiation, with positive logarithm arguments, yields
\begin{align}
 c'(r)&=\tfrac12\log\frac{1+r}{1-r}-\frac{r}{(1+r)^2},\label{cp}\\
 c''(r)&=\frac1{1-r^2}-\frac{1-r}{(1+r)^3}
       =\frac{4r}{(1-r)(1+r)^3}.\label{cpp}
\end{align}
The formulas extend to $r=0$ with $c'(0)=0$. Equation \eqref{cpp} is strictly positive for $0<r<1$. By integration, or the mean value theorem applied first to $c'$ and then to $c$, $c$ is strictly increasing on $[0,1)$. This is not a claim that $c$ is positive: $c(0)=-1/2$.

\begin{lemma}\label{negative}
For every $0<l<1$,
\begin{equation}
 2\sum_{i=1}^{\infty}2^{-i}[c(l^i)-c(l)]<0.\label{negseries}
\end{equation}
\end{lemma}
\begin{proof}
All arguments $l^i$ lie in $[0,l]$, so $c$ is bounded there and the series is absolutely convergent. Since $l^i\leq l$, every term is nonpositive. The $i=2$ term is strictly negative by $l^2<l$ and strict monotonicity. The sum is bounded above by this strictly negative term, proving \eqref{negseries}.
\end{proof}

\section{Uniform differentiation and promotion to two dimensions}
We now justify using the coefficient calculation as an actual curvature calculation for the infinite series. Fix $b<1$ with $b\geq0$, and $0<\delta<1$. For every $r\in[0,b]$, $|t|\leq\delta$, both arguments in \eqref{H} lie in
\begin{equation}
 K=[A,B],\qquad A=\frac{(1-b)(1-\delta)}2>0,
 \quad B=\frac{1+\delta}2<1.\label{compactK}
\end{equation}
Choose $M_j=\max_{y\in K}|h^{(j)}(y)|$ for $j=0,1,2$, which are finite by continuity on this compact interval. Because $0\leq x\leq1/2$, \eqref{first}--\eqref{second} give, uniformly on the entire rectangle,
\begin{equation}
 |H|\leq\log2,\quad
 |H_t|\leq M_0+\tfrac12(1+\delta)M_1=:C_1,\quad
 |H_{tt}|\leq M_1+\tfrac14(1+\delta)M_2=:C_2.\label{bounds}
\end{equation}
The slightly larger constants displayed suffice even without optimizing the sum of the weights. In particular these bounds include $r=0$, the limit of the gap powers; a rectangle with $r$ bounded away from zero would not suffice.

For $0\leq l\leq b$ and every $i\geq1$, $0\leq l^i\leq b$. Thus $2^{-i}\log2$, $2^{-i}C_1$, $2^{-i}C_2$ are summable majorants for the series in \eqref{F} and its first two transverse derivative series. Here is an explicit justification of differentiation. For fixed $l$, finite partial sums $G_m(t)=\sum_{i=1}^m2^{-i}H(l^i,t)$ satisfy
\[
 G_m(t)-G_m(0)=\int_0^t G_m'(s)\,ds.
\]
Uniform convergence of both sides' series allows passage to the limit, uniformly for $|t|\leq\delta$. The derivative series has a continuous sum by uniform convergence, so the fundamental theorem of calculus identifies it as the derivative of the function series on $(-\delta,\delta)$. Apply the same argument to $G_m'$ using uniform convergence of $G_m''$. Subtracting the single term $H(l,t)$ therefore gives
\begin{equation}
 F_{tt}(l,t)=\sum_{i=1}^{\infty}2^{-i}H_{tt}(l^i,t)-H_{tt}(l,t).\label{Ftt}
\end{equation}
Each summand in \eqref{Ftt} is jointly continuous in $(l,t)$ on $[0,b]\times[-\delta,\delta]$, including $l=0$, and \eqref{bounds} proves uniform convergence there. Thus $F_{tt}$ is jointly continuous. No derivative in $l$ is needed.

Combining \eqref{coefficient}, \eqref{Ftt}, $\sum_{i\geq1}2^{-i}=1$ and Lemma~\ref{negative} proves
\begin{equation}
 F_{tt}(l,0)=2\sum_{i=1}^{\infty}2^{-i}[c(l^i)-c(l)]<0
 \quad(0<l<1).\label{curvature}
\end{equation}
Choose $l_0$ as in \eqref{maxdiag}, and choose $b\in(l_0,1)$, say $b=(1+l_0)/2$, and initially $\delta=1/2$. Joint continuity and \eqref{curvature} give $\epsilon>0$, $\eta>0$, with
\[
 (l_0-\epsilon,l_0+\epsilon)\subset(0,b),\quad \eta<1/2,
 \quad F_{tt}(l,t)<0\quad\text{if }|l-l_0|<\epsilon,\ |t|<\eta.
\]
Interchanging the two labels exchanges the two terms of \eqref{H}, so $H(r,-t)=H(r,t)$ and $F(l,-t)=F(l,t)$. For fixed $l$ in this rectangle, the twice differentiable transverse function is concave. Applying its midpoint inequality to $t,-t$ yields $F(l,t)\leq F(l,0)$. Alternatively evenness gives $F_t(l,0)=0$ and integration of the negative second derivative gives the same inequality, strictly for $t\ne0$. Equation \eqref{maxdiag} then gives
\begin{equation}
 F(l,t)\leq F(l,0)\leq F(l_0,0)\qquad
 (|l-l_0|<\epsilon,\ |t|<\eta).\label{2d}
\end{equation}
This establishes a full two-dimensional local maximum. It allows a flat or nonunique diagonal maximizing set.

\section{Transfer to every nearby stationary one-step law}
Let $\mu_* =\mu_{l_0,0}$. For a probability law $\nu$ write
\[
 m_a(\nu)=\nu[x_0=a],\qquad d_a(\nu)=\nu[x_0=a,x_1=a]
 \quad(a=0,1).
\]
These cylinders are clopen in the compact product space $X$, so their indicator functions are continuous bounded functions. Consequently $m_a,d_a$ are continuous for weak convergence of probability laws. At $\mu_*$ both $m_a$ equal $1/2$. On the open set $m_0m_1>0$ the recovered diagonals
\[
 R(\nu)=\left(\frac{d_0(\nu)}{m_0(\nu)},
                 \frac{d_1(\nu)}{m_1(\nu)}\right)
\]
are therefore continuous. At $\mu_*$ they equal $(u_*,v_*)=((1+l_0)/2,(1+l_0)/2)$.

Let $V$ be an open neighborhood of $(u_*,v_*)$ small enough that every point is interior to the unit square, has $2-u-v>0$, and has inverse chart in the rectangle used in \eqref{2d}. Such a $V$ exists by \eqref{inverse}, continuity and $0<l_0<1$. Set
\[
 U=\{\nu:m_0(\nu)>0,\ m_1(\nu)>0,\ R(\nu)\in V\}.
\]
This is an open weak neighborhood of $\mu_*$. Suppose $\nu\in U\cap\mathcal M_1$, and choose \eqref{block} data $p,P$ for it, with $P_{00}=u$, $P_{11}=v$. The length-one and length-two cylinder identities give
$m_a(\nu)=p_a$ and $d_a(\nu)=p_aP_{aa}$. Positive $m_a$ implies $R(\nu)=(u,v)$, independently of the chosen presentation of the law. Since $(u,v)\in V$, the stationary vector is uniquely \eqref{stationary}. The inverse parameters $(l,t)$ satisfy $0<l<1$, $|t|<1$ and $P=P(l,t)$, $p=\pi_t$. Hence \eqref{block} agrees with that of $\mu_{l,t}$ on every translated interval.

To see that this gives equality of laws without an extra assumption, embed any finite set of integer indices in a finite interval and sum the interval atoms to recover its cylinder probabilities. Thus the laws agree on all finite-coordinate cylinders. This generating pi-system includes $X$, and the laws are finite measures of equal total mass; measure uniqueness gives $\nu=\mu_{l,t}$. Proposition~\ref{series} proves actual finite-limit convergence for this $\nu$, and \eqref{2d} proves $\Int(\nu)\leq\Int(\mu_*)$. Full support and invariance of $\mu_*$ were proved in Section 2. This completes Theorem~\ref{main}.

The comparison class in this proof really is all of $\mathcal M_1$. For example mixtures of the two constant sequences ($P=I$) and the stationary deterministic alternating law belong to it. Such degeneracies are excluded from $U$ by proved positive-mass/interior-recovery conditions, not by an extra hypothesis in the definition of $\mathcal M_1$. No assertion about convergence at remote boundary laws is necessary for the local assertion.

\section{Formal certificate, reproduction and trust}
The dossier preserves every supplied Lean byte and the original
\texttt{proof-packet.zip}. The endpoint is
\begin{center}\small\texttt{MathDB333569.fullySupported\_stationary\_oneStep\_localMaximum}.\end{center}
It quantifies over actual probability measures on $\mathbb Z\to\operatorname{Fin}2$, asserts full open support, invariance under every integer translation, the full cylinder-law predicate, convergence at the witness and all nearby comparison laws, and relative weak local maximality. Although the definition \texttt{actualIntricacy} uses \texttt{limUnder}, those explicit convergence clauses ensure that the compared values here are genuine limits. A mere equality between two unproved \texttt{limUnder} expressions would not suffice.

The 21 mathematical modules supply respectively the matrix chart and algebra, coefficient calculus, diagonal compactness, coefficient-series sign, finite word/selected entropies, actual two-sided measure construction and stationarity, exact fair-subset counts and limit, actual transverse derivatives, uniform twice differentiation and continuity, parameter maximum, inverse chart and weak-topology recovery. The frozen source mapping and independent semantic review identify these bridges by declaration name.

Lean is pinned to 4.34.1 and mathlib to revision
\begin{center}\small\texttt{d13f23b723b8a846827a245b89c10fc7d3f11612}.\end{center}
The dossier validator safely extracts the original packet with path, symlink, size and compiled-artifact rejection; fresh local compilation uses a sanitized environment, two Lean threads and at most two CPUs, a 4096 MB Lean memory cap, a 180 CPU-second limit and a 240-second wall limit per module. It uses only the explicitly pinned external dependency paths plus its new extraction directory. Logs and resource receipts accompany the dossier.

The original author regex inventories 238 source-written named declarations. The augmented validator instead enumerates the rebuilt Lean environment, filters by originating mathematical module, records every declaration's kind and unsafe flag, and calls \texttt{Lean.collectAxioms} on each. This includes 206 additional generated declarations, for 444 total. It checks coverage of all 238 names, rejects local axioms and unsafe declarations, and allows only \texttt{propext}, \texttt{Classical.choice}, \texttt{Quot.sound} transitively. The independent review already found all 444 compliant; the dossier provides its own fresh replay and the same complete method rather than treating a text regex as comprehensive.

Prebuilt upstream Lean/mathlib/package oleans and their exported axiom summaries remain trusted. Package metadata and source pins do not attest source-to-olean identity. Neither this dossier nor the prior review rebuilds the entire upstream library or uses an independent kernel. The local sources are freshly checked against that explicit cache. Paper review, root final replay and root acceptance are separate; this draft awards no campaign credit.

\section{Primary comparisons and limits on novelty claims}
\raggedright
The normalization, chain rule, Markov powers, geometric series and compactness infrastructure are standard. In the inspected sources the proposed substantive addition is the explicit increasing transverse coefficient \eqref{coefficient}--\eqref{cpp}, giving strict curvature \eqref{curvature} throughout the positively correlated symmetric family and closing the two-dimensional local-existence argument. Full Lean completion is not evidence that this addition is new to the literature.

The qualified source review read Petersen--Wilson's Proposition 5.10 and its proof (5.28)--(5.31): conditional-entropy telescoping, the Markov property and geometric return-time weights yield the entropy series. Section 6 motivates conjectures by numerical plots using 20 terms. Its zero-intricacy calculation (6.8)--(6.10) concerns the independent line $u=1-v$ ($l=0$), not transverse stability on $u=v$. Proposition 5.5 obtains an Asc-maximizer among invariant laws by affinity, upper semicontinuity and compactness; it supplies neither an Int-maximizer on the Markov locus nor the curvature/support conclusion here. Petersen's Sections 4.5--4.6 likewise provide the formulas and numerical symmetric candidate, not the displayed sign proof. These are comparisons with the inspected cached versions, not claims about every later version.

The older review read selected primary proof bodies of Buzzi--Zambotti arXiv:0909.2120v2 \cite{BZ}: finite-law definitions, Proposition 3.3, the random sparse-support construction and proof of Proposition 4.3, and Theorem 5.1. Their laws on finite strings may depend on string length; the occupancy and entropy estimates do not produce consistent stationary one-step Markov marginals or this local two-parameter curvature result. The symmetric entropy-profile envelope is not concavity of $F(l,t)$.

The same review compared the definitions and complete proofs of Huang--Xiao arXiv:2509.20738v1, Section 5, Theorems 5.1--5.2 \cite{HX}. The first uses the system-level supremum over covers/partitions and the classical variational principle; the second obtains a fixed-cover Asc-maximizer and an ergodic component by affinity. Neither inspected mechanism supplies subtraction-of-entropy optimization on the one-step locus. Their use of ``local'' for a fixed cover is a different quantifier from a local extremum in Markov-law space. No complete subsumer was found in those inspected bodies; this is a bounded statement, not global clearance.

Explicit limits remain: Wilson's 2015 dissertation body was not located; the related Buzzi--Zambotti 2012 symmetry-breaking proof body was unread; Li--Tu and other citing/pressure/random-bundle bodies were only partially covered. The Petersen--Wilson publisher full text returned public HTTP 403 and was not retried or bypassed. The Huang--Xiao publisher landing page was accessible, but its full published proof was not compared. All earlier versions and all citing literature were not audited. The MathDB/Palomar public access denials were respected. The reused OpenAI registry receipt dated 1 November 2026 authenticates only a recent head response, not a current repository-wide proof-body comparison. Empty solution arrays and open labels provide no novelty evidence. A prior equivalent binary-channel inequality or transverse-concavity theorem with a complete substitution could change the significance judgment. The dossier retains the qualified review and its receipts; no worldwide priority claim or new clearance is made here.

\begin{thebibliography}{9}
\bibitem{P} K. Petersen, \emph{Measuring Complexity in Cantor Dynamics}, arXiv:1607.02425v1, 8 July 2016, Sections 4.5--4.6, Conjecture 4.6.8. \url{https://arxiv.org/abs/1607.02425v1}.
\bibitem{PW} K. Petersen and B. Wilson, \emph{Dynamical intricacy and average sample complexity}, arXiv:1512.01143v6, 18 November 2016, Definition 5.1, Propositions 5.5 and 5.10, Section 6, Conjecture 6.3. \url{https://arxiv.org/abs/1512.01143v6}.
\bibitem{BZ} J. Buzzi and L. Zambotti, \emph{Approximate maximizers of intricacy functionals}, arXiv:0909.2120v2, selected statements and proofs as specified above. \url{https://arxiv.org/abs/0909.2120v2}.
\bibitem{HX} J. Huang and Z. Xiao, \emph{Local intricacy and average sample complexity for amenable group actions}, arXiv:2509.20738v1, Section 5, Theorems 5.1--5.2 and their proofs, bounded comparison described above. \url{https://arxiv.org/abs/2509.20738v1}.
\end{thebibliography}
\end{document}
